% !TeX root = ../教材题库_全部条目_转写工作稿.tex
% Source: XJ-B1-C05-EX-P167-Y0891
% 第5章 / 5.2 任意角的三角函数 / 教材 p.167 / PDF 第174页
% OCR未校对：公式与文字均可能有误；未完成数学TeX转写；禁止直接发给学生
% 原题及相邻原解截图（核对用）：../题目原图/XJ-B1-C05-EX-P167-Y0891-01.webp
% 本文件只在转写完成后作为题干使用；OCR 条目尚不是可靠题干。
\begin{bankitem}{5.2 例题 例题4}{第5章 / 5.2 任意角的三角函数 / 教材 p.167 / PDF 第174页}{XJ-B1-C05-EX-P167-Y0891}{OCR未校对：公式与文字均可能有误；未完成数学TeX转写；禁止直接发给学生}
\begin{ocrdraft}
[教材p.167 / PDF第174页]
\par
0互利用正弦线、余弦线、正切线研究各象限角的三角函数值的符号．
\par
解观察图5．2一6中各象限角的三角函数线，可知：
\par
第一、二象限角的正弦线DP为正向，正弦值为正；第三、四象限角的正弦线
\par
DP为负向，正弦值为负．
\par
第一、四象限角的余弦线OD为正向，余弦值为正；第二
\par
三象限角的余弦线
\par
OD为负向，余弦值为负．
\par
第一、三象限角的正切线AT为正向，正切值为正；第二、四象限角的正切线
\par
，AT为负向，正切值为负．
\end{ocrdraft}
\end{bankitem}
